% Source: https://artofproblemsolving.com/community/c3h1984745_all_digits_of_a_number
% Stable identity: aops-topic-1984745-q01
The last two digits are the remainder of $3^{678}$ modulo $100$. Since $3^{10}=59049$, we have $3^{10}\equiv49$ and $3^{20}\equiv49^2=2401\equiv1\pmod{100}$. As $678=20\cdot33+18$,
\[
3^{678}\equiv3^{18}=3^{10}\cdot3^{8}\equiv49\cdot6561\equiv49\cdot61=2989\equiv89\pmod{100}.
\]
So the last two digits are $\boxed{89}$. In the same way, the last $k$ digits of any power are its remainder modulo $10^k$; finding leading digits needs logarithms instead.
